%-----------------------------------------------------------------------------------------------
\makeatletter
\immediate\write18{datelog > \jobname.info} % site script for $(date '+%Y-%m-%d %Hh%Mm%Ss')
\makeatother
%-----------------------------------------------------------------------------------------------
%-----------------------------------------------------------------------------------------------
\usetheme{Copenhagen}
\usepackage{beamercolorthemeCNThermSci2}
\usefonttheme{serif}
%-----------------------------------------------------------------------------------------------

%-----------------------------------------------------------------------------------------------
%-----------------------------------------------------------------------------------------------
\usetheme{Copenhagen}
\usepackage{beamercolorthemeUTF2}
\usefonttheme{serif}
%-----------------------------------------------------------------------------------------------
\usepackage[utf8]{inputenc}
\usepackage[greek,french,english,brazil]{babel} % last becomes the active one
\usepackage{pslatex}
\usepackage{amssymb,amsmath}
\usepackage{soul}
\usepackage[squaren,Gray,cdot]{SIunits}
\usepackage[nice]{nicefrac}
\usepackage{tikz}
\usepackage{amscd}
\usepackage{stmaryrd}
\usepackage{scalerel}
\usepackage{xspace}
%-----------------------------------------------------------------------------------------------


%-----------------------------------------------------------------------------------------------
%-----------------------------------------------------------------------------------------------
% Mathematical
%-----------------------------------------------------------------------------------------------
\newcommand{\vet}[1]{\underline{{#1}}}
\newcommand{\mat}[1]{\underline{\underline{{#1}}}}
\newcommand{\cub}[1]{\underline{\underline{\underline{{#1}}}}}
\newcommand{\eqdef}{{\ensuremath\stackrel{\text{\tiny def}}{=}}}
%-----------------------------------------------------------------------------------------------
% Linguistic
%-----------------------------------------------------------------------------------------------
\newcommand{\GRtxt}[1]{\begin{otherlanguage}{greek}{{#1}}\end{otherlanguage}}
\newcommand{\FRtxt}[1]{\begin{otherlanguage}{french}{{#1}}\end{otherlanguage}}
%-----------------------------------------------------------------------------------------------
% Presentation
%-----------------------------------------------------------------------------------------------
\newcommand{\BkgImgH}[1]{% Places an image centered on the slide background filling the height
    \usebackgroundtemplate{\parbox{\paperwidth}{%
        \vspace*{1sp}\centering\includegraphics[height=\paperheight]{{#1}}
}}}
\newcommand{\BkgImgW}[1]{% Places an image centered on the slide background filling the width
    \usebackgroundtemplate{\parbox{\paperwidth}{%
        \vspace*{1sp}\centering\includegraphics[width=\paperwidth]{{#1}}
}}}
\newcommand{\ArtEndH}[3]{% Transitions to plain image (last) slide: #1:prefix #2,#3:extensions
    \BkgImgH{root/../art/#1.#2}
    \frame<handout:0>[plain]{%
        \transdissolve\vspace*{72mm}\color{white}\scriptsize\bf\input{root/../art/#1.#3}}
    \usebackgroundtemplate{\mbox{~}}
}
\newcommand{\ArtEndW}[3]{% Transitions to plain image (last) slide: #1:prefix #2,#3:extensions
    \BkgImgW{root/../art/#1.#2}
    \frame<handout:0>[plain]{%
        \transdissolve\vspace*{72mm}\color{white}\scriptsize\bf\input{root/../art/#1.#3}}
    \usebackgroundtemplate{\mbox{~}}
}
\newcommand{\ImgColW}[3]{% Inserts a full-width image in a column
    \includegraphics[width=\columnwidth]{root/../art/#1.#2}\\[-0.5\baselineskip]
    \parbox{\columnwidth}{\tiny\hfill\scalebox{0.85}{\input{root/../art/#1.#3}}}
}
\newcommand{\txtpic}[1]{%
    \fcolorbox{lightgray}{white!90!black}{{#1}} 
}
%-----------------------------------------------------------------------------------------------


%-----------------------------------------------------------------------------------------------
\title{A.08.01 -- Misturas Gás-Vapor e Condicionamento de Ar}
\subtitle{Ar Seco e Atmosférico e Medidas de Umidade}
\author{Prof.~C.~Naaktgeboren, PhD}
\date{{\scriptsize\tt%
    \includegraphics[height=6.0mm]{cc/by-nc-nd-88x31.pdf}\\[\smallskipamount]
    https://github.com/CNThermSci/ApplThermSci\\
    Compiled on \input{\jobname.info}
}}
%-----------------------------------------------------------------------------------------------
\begin{document}
%-----------------------------------------------------------------------------------------------
\logo{%
    \parbox{158mm}{% There's a 1mm gap on each side of the 160mm x 90mm slide logo line
        \mode<beamer>{
            \includegraphics[height=6.0mm]{root/00-res/UTFPR/UTFPR-logo-D.pdf}\hfill%
            \includegraphics[height=9.0mm]{root/00-res/logo/CNThermSci-logo-A.pdf}%
        }
        \mode<handout>{
            \includegraphics[height=6.0mm]{root/00-res/UTFPR/UTFPR-logo-W.pdf}\hfill%
            \includegraphics[height=9.0mm]{root/00-res/logo/CNThermSci-logo-W.pdf}%
        }
    }
} % The (delineated, alpha), or washed-out logos
%-----------------------------------------------------------------------------------------------
\frame{\titlepage}
%-----------------------------------------------------------------------------------------------

%-----------------------------------------------------------------------------------------------
\frame{\tableofcontents}
%-----------------------------------------------------------------------------------------------

%-----------------------------------------------------------------------------------------------
\section{Definições}
%-----------------------------------------------------------------------------------------------

%-----------------------------------------------------------------------------------------------
\subsection{Ar Seco e Ar Atmosférico}
%-----------------------------------------------------------------------------------------------

    % !j 96 -i8
    %-------------------------------------------------------------------------------------------
    \newcommand{\THISFRAMECONTENTS}{
        \begin{itemize}
            \item<1-> \alert{Ar atmosférico} é uma mistura de \textcolor{black}{\ce{N2},
                \ce{O2},} outros gases e \alert{vapor d'água};
            \item<2-> A quantidade de outros gases e vapor d'água é pequena;
            \item<3-> Ar \alert{isento} de vapor d'água é chamado de \alert{ar seco};
            \item<4-> Ar atmosférico é modelado como uma mistura de (i)~\alert{ar seco} e
                (ii)~\alert{vapor d'água};
            \item<5-> Já que a composição do ar seco em tal modelo é estável;
            \item<6-> E as interações \alert{energéticas} do \alert{vapor d'água} serem
                \alert{importantes} e \alert{distintas}; \\[\smallskipamount]
            \begin{itemize}
                \item<7-> Maior \alert{calor específico}: $c_{P,v}$ é \alert{$81\%$ maior} que
                    $c_{P,a}$ (base mássica): \\[\smallskipamount]
                \item<8-> $c_{P,a} \simeq \unit{1,005}{\kilo\joule\per\kilogram}$
                    $c_{P,v} \simeq \unit{1,82}{\kilo\joule\per\kilogram}$; \\[\smallskipamount]
                \item<9-> \alert{Calor latente}: condensação e evaporação.
            \end{itemize}
        \end{itemize}
    }
    % Presentation
    {\usebackgroundtemplate{\parbox{\paperwidth}{%
        \vspace*{1sp}\centering
        \includegraphics[height=\paperheight]{root/../art/pexels-neosiam-601798.jpg}
    }}
    \begin{frame}<handout:0>{Ar Seco e Ar Atmosférico}\vspace*{-0em}
        \THISFRAMECONTENTS
    \end{frame}}
    % Handout
    {\usebackgroundtemplate{\parbox{\paperwidth}{%
        \vspace*{1sp}\centering
        \includegraphics[height=\paperheight]{root/../art/wo-pexels-neosiam-601798.jpg}
    }}
    \begin{frame}<beamer:0>{Ar Seco e Ar Atmosférico}\vspace*{-0em}
        \THISFRAMECONTENTS
    \end{frame}}
    %-------------------------------------------------------------------------------------------

    %-------------------------------------------------------------------------------------------
    \begin{frame}{Ar Seco e Ar Atmosférico --- Aproximações}\vspace*{-0em}
        Nas \alert{estreitas faixas de temperatura} pertinentes ao \alert{condicionamento de
        ar}, as seguintes aproximações são aceitáveis (na quais $\mathsf{T}$ estão em \celsius):
        \begin{align*}
            \alert{h_{ar,seco}} &
            \alert{
                \simeq c_{P,a}\mathsf{T}
                = (\unit{1,005}{\kilo\joule\per\kilogram\celsius})\mathsf{T}
            }
            \\[\medskipamount]
            \alert{\Delta h_{ar,seco}} &
            \alert{
                \simeq c_{P,a}\Delta\mathsf{T}
                = (\unit{1,005}{\kilo\joule\per\kilogram\celsius})\Delta\mathsf{T}
            }
            \\[\medskipamount]
            \alert{h_{v}} &
            \alert{
                \simeq \unit{2500,9}{\kilo\joule\per\kilogram}
                + (\unit{1,82}{\kilo\joule\per\kilogram\celsius})\mathsf{T}
            }
        \end{align*}
    \end{frame}
    %-------------------------------------------------------------------------------------------

%-----------------------------------------------------------------------------------------------
\subsection{Pressão de Vapor}
%-----------------------------------------------------------------------------------------------

    %-------------------------------------------------------------------------------------------
    \begin{frame}{Pressão de Vapor}\vspace*{-0em}
        Para \alert{comportamento $P-V-T$ ideal} da mistura, tem-se a \alert{pressão
        componente}:
        \begin{align*}
            \alert{\frac{P_{i}(T_m, V_m)}{P_m}} &
            \alert{=\frac{N_{i}}{N_m} = y_{i}}
            \\[\medskipamount]
            \alert{\frac{P_i}{P_m}} &
            \alert{=\frac{V_i}{V_m} = \frac{N_i}{N_m} = y_i},
        \end{align*}
        assim como a \alert{pressão parcial}, \alert{$y_iP_m$}, a qual, aplicada ao vapor do ar
        atmosférico é também chamada de \alert{pressão de vapor}, \alert{$P_v = y_vP$}, com:
        \begin{equation*}
            \alert{
                P = P_a + P_v
            }.
        \end{equation*}
    \end{frame}
    %-------------------------------------------------------------------------------------------

%-----------------------------------------------------------------------------------------------
\subsection{Umidade Específica e Relativa do Ar}
%-----------------------------------------------------------------------------------------------

    %-------------------------------------------------------------------------------------------
    \begin{frame}{Umidade Específica do Ar}\vspace*{-0em}
        A razão entre as massas de vapor d'água pela de ar seco no ar atmosférico é chamada de
        \alert{umidade absoluta} ou \alert{específica}\footnote{Esta razão também é chamada de
        \alert{relação de umidade}.}:
        \begin{align*}
            \alert{\omega} &
            \alert{\equiv\frac{m_v}{m_a}}
            = \frac{P_vV/R_vT}{P_aV/R_aT} = \frac{R_a}{R_v}\frac{P_v}{P_a}
            \\[\medskipamount]
            \alert{\omega} &
            \alert{= \frac{M_v}{M_a}\frac{P_v}{P_a}}
            \simeq
            \frac{\unit{18,015}{\kilogram\per\kilo\mole}}{\unit{28,97}{\kilogram\per\kilo\mole}}
            \frac{P_v}{P_a}
            \alert{\simeq 0,622\frac{P_v}{P_a}},
        \end{align*}
        ou, substituindo $P$ em $P_a$,
        \begin{equation*}
            \alert{
                \omega = \frac{0,622P_v}{P-P_v}\qquad
                \mbox{($\kilogram$ de vap.$/\kilogram$ de ar seco)}
            }.
        \end{equation*}
    \end{frame}
    %-------------------------------------------------------------------------------------------

    %-------------------------------------------------------------------------------------------
    \begin{frame}{Umidade Relativa do Ar}\vspace*{-0em}
        \begin{itemize}
            \item<1-> A máxima umidade do ar a uma certa $T$ ocorre quando o vapor é
                \alert{saturado};
            \item<1-> A saber, quando \alert{$P_v = P_g = P_{\mathrm{sat}\,@\,T}$};
            \item<1-> Na aplicação de \alert{conforto térmico}, a \alert{unidade relativa},
                \alert{$\phi$} é mais significativa:
        \end{itemize}
        \begin{align*}
            \alert{\phi} &
            \alert{\equiv \frac{m_v}{m_g}}
            =\frac{P_vV/R_vT}{P_gV/R_vT} = \alert{\frac{P_v}{P_g}}
            \qquad\rightharpoondown\qquad
            \alert{P_v = \phi P_g}
            \\[\medskipamount]
            \alert{\phi} &
            \alert{=\frac{\omega P}{(0,622+\omega)P_g}},
            \qquad\mbox{e}\qquad
            \alert{\omega=\frac{0,622\phi P_g}{P-\phi P_g}}.
        \end{align*}
    \end{frame}
    %-------------------------------------------------------------------------------------------

    %-------------------------------------------------------------------------------------------
    \begin{frame}{Entalpia do Ar Atmosférico}\vspace*{-0em}
        \begin{align*}
            H &= H_a + H_v = m_ah_a + m_vh_v
            \qquad\mbox{(propr.~ext.)}
            \\[\medskipamount]
            h &= \frac{H}{m_a} = h_a + \frac{m_v}{m_a}h_v
            \qquad\rightharpoondown
            \\[\medskipamount]
            \alert{h} &
            = h_a + \omega h_v
            \alert{\simeq h_a + \omega h_g}
            \qquad\mbox{(gás ideal: $h\!:\!h(T)$)}.
        \end{align*}
    \end{frame}
    %-------------------------------------------------------------------------------------------

%-----------------------------------------------------------------------------------------------
\section{Tópicos de Leitura}
%-----------------------------------------------------------------------------------------------

    %------------------------------------------------------------------------------------------
    \begin{frame}[allowframebreaks]{Tópicos de Leitura}
        \begin{thebibliography}{Çengel, Y.~A., 2013}
            \bibitem[Çengel, Y.~A., 2013]{2013-CengelYA+BolesMA-AMGH}
                Çengel, Y.~A. e Boles, M.~A.
                \newblock{%
                    {\em Termodinâmica $7^\mathrm{a}\!$ Edição\/}.
                    \alert{Seções~14-1 a 14-2.}
                }
                \newblock{\footnotesize AMGH. Porto Alegre. ISBN 978-85-8055-200-3.}
        \end{thebibliography}
    \end{frame}
    %------------------------------------------------------------------------------------------

    % Finishes with stunning image, with credit
    \ArtEndW{pexels-pixabay-268917}{jpg}{txt}

%-----------------------------------------------------------------------------------------------
\end{document}
%-----------------------------------------------------------------------------------------------

